\subsection{张量积的结合性}

命题8. 设 A、B 为两个环，E 为右 A-模，F 为 (A, B)-双模，G 为左 B-模。则 E \(\otimes_{\mathbf{A}} \mathbf{F}\) 是右 B-模，F \(\otimes_{\mathbf{B}} \mathbf{G}\) 是左 A-模，并且存在且仅存在一个 Z-线性映射

\[
\phi : (\mathrm{E} \otimes_ {\mathrm{A}} \mathrm{F}) \otimes_ {\mathrm{B}} \mathrm{G} \rightarrow \mathrm{E} \otimes_ {\mathrm{A}} (\mathrm{F} \otimes_ {\mathrm{B}} \mathrm{G})
\]

，使得 \(\phi((x \otimes y) \otimes z) = x \otimes (y \otimes z)\) （对 \(x \in \mathbf{E}\) , \(y \in \mathbf{F}\) , \(z \in \mathbf{G}\) ）；此外，这个 \(\mathbf{Z}\) -线性映射是双射的（张量积的“结合性”）。

\(\mathbf{E} \otimes_{\mathbf{A}} \mathbf{F}\) 上的右 B-模结构和 \(\mathbf{F} \otimes_{\mathbf{B}} \mathbf{G}\) 上的左 A-模结构已在第4号中定义。 \(\phi\) 的唯一性是显然的，因为元素 \((x \otimes y) \otimes z\) 生成 Z-模 \((\mathbf{E} \otimes_{\mathbf{A}} \mathbf{F}) \otimes_{\mathbf{B}} \mathbf{G}\) 。为了证明 \(\phi\) 的存在性，我们注意到，对所有 \(z \in \mathbf{G}\) , \(h_z: y \mapsto y \otimes z\) 是从左 A-模 F 到左 A-模 \(\mathbf{F} \otimes_{\mathbf{B}} \mathbf{G}\) 的 A-线性映射。我们记 \(g_z = l_{\mathbf{E}} \otimes h_z\) ，它因此是 \(\mathbf{E} \otimes_{\mathbf{A}} \mathbf{F}\) 到 \(\mathbf{E} \otimes_{\mathbf{A}} (\mathbf{F} \otimes_{\mathbf{B}} \mathbf{G})\) 的一个 Z-线性映射，并考虑映射 \((t, z) \mapsto g_z(t)\) （从 \((\mathbf{E} \otimes_{\mathbf{A}} \mathbf{F} \times \mathbf{G}\) 到 \(\mathbf{E} \otimes_{\mathbf{A}} (\mathbf{F} \otimes_{\mathbf{B}} \mathbf{G})\) ）；由于 \(h_{z+z'} = h_z + h_{z'}\) （对 \(z \in \mathbf{G}\) , \(z' \in \mathbf{G}\) ），显然上述映射是 Z-双线性的。进一步，我们证明对所有 \(\mu \in \mathbf{B}\) , \(g_{\mu z}(t) = g_z(t\mu)\) 成立；显然只需对 \(t = x \otimes y\) （其中 \(x \in \mathbf{E}\) 和 \(y \in \mathbf{F}\) ）证明即可；现在

\[
g _ {\mu z} (x \otimes y) = x \otimes (y \otimes \mu z)
\]

\[
g _ {z} ((x \otimes y) \mu) = g _ {z} (x \otimes y \mu) = x \otimes (y \mu \otimes z).
\]

命题1（第 1 号）随即证明了 Z-线性映射

\[
\phi : (\mathrm{E} \otimes_ {\mathrm{A}} \mathrm{F}) \otimes_ {\mathrm{B}} \mathrm{G} \rightarrow \mathrm{E} \otimes_ {\mathrm{A}} (\mathrm{F} \otimes_ {\mathrm{B}} \mathrm{G})
\]

，使得 \(\phi(t \otimes z) = g_z(t)\) ，因此 \(\phi((x \otimes y) \otimes z) = x \otimes (y \otimes z)\) 的存在性。类似地，一个 Z-线性映射

\[
\psi : \mathbf {E} \otimes_ {\mathbf {A}} (\mathbf {F} \otimes_ {\mathbf {B}} \mathbf {G}) \rightarrow (\mathbf {E} \otimes_ {\mathbf {A}} \mathbf {F}) \otimes_ {\mathbf {B}} \mathbf {G}
\]

定义为使 \(\psi(x \otimes (y \otimes z)) = (x \otimes y) \otimes z\) ，并且显然 \(\psi \circ \phi\) 和 \(\phi \circ \psi\) 分别是 \((\mathbf{E} \otimes_{\mathbf{A}} \mathbf{F}) \otimes_{\mathbf{B}} \mathbf{G}\) 和 \(\mathbf{E} \otimes_{\mathbf{A}} (\mathbf{F} \otimes_{\mathbf{B}} \mathbf{G})\) 的恒等映射，因为它们在这些 \(\mathbf{Z}\) -模的生成系上退化为恒等映射。

显然，如果 E 是一个 \(((C_{i}^{\prime}); A, (D_{j}^{\prime}))\) -多重模，F 是一个 \((A, (C_{h}^{\prime\prime}); B, (D_{k}^{\prime\prime}))\) -多重模，G 是一个 \((B, (C_{l}^{\prime\prime}); (D_{m}^{\prime\prime}))\) -多重模，则命题 8 中定义的典范同构是一个 \(((C_{i}^{\prime}), (C_{h}^{\prime\prime}), (C_{l}^{\prime\prime}); (D_{j}^{\prime}), (D_{k}^{\prime\prime}), (D_{m}^{\prime\prime}))\) -多重模同构。特别地，若 C 是交换环且 E、F、G 是三个 C-模，则存在典范的 C-模同构

\[
(\mathbf {E} \otimes_ {\mathbf {c}} \mathbf {F}) \otimes_ {\mathbf {c}} \mathbf {G} \rightarrow \mathbf {E} \otimes_ {\mathbf {c}} (\mathbf {F} \otimes_ {\mathbf {c}} \mathbf {G}).
\]

我们将在下面看到，在一定条件下，张量积的定义可以推广到一族多重模，这将在命题 8 的假设下特别给出一个 Z-模 \(E \otimes_{A} F \otimes_{B} G\) ，它与 \((E \otimes_{A} F) \otimes_{B} G\) 和 \(E \otimes_{A} (F \otimes_{B} G)\) 两个 Z-模中的每一个都典范同构，并且后者将与它等同。

